\documentclass[fleqn]{article}
\usepackage{graphicx} % Required for inserting images

\usepackage{amsmath}

\title{CS215 Assignment-1}
\author{Malay Kedia}
\date{August 2024}

\begin{document}
\maketitle

\section{Solution Q1}

Let $w_{P, k}$ represent the event that player $P$ gets a prime number on $k^{\text {th}}$ throw, and hence gets a win.
\[
w_{P, k} = P(2,3,5)=\frac{3}{6}=\frac{1}{2} \text{ for both players A and B}
\]
Let $X_r$ and $Y_r$ denote the random variables representing no of wins by $A$ and $B$ respectively in first $r$ rolls.  Now,
\[
P(X_{n+1}>Y_n)=P(X_n>Y_n)\cdot(1)+P(X_n=Y_n) \cdot P(w_{A, n+1})
\]
as either A has more wins in first n rolls itself, or A gets the last win in $(n+1)^{th}$ roll.\\
Now,
\begin{align*}
P(X_n = Y_n = r) &= P(X_n = r) \cdot P(Y_n = r) = P(X_n = r)^2 \\
&= \left(\binom{n}{r} w^r (1-w)^{n-r}\right)^2 = \left(\binom{n}{r} \left(\frac{1}{2}\right)^r \left(\frac{1}{2}\right)^{n-r}\right)^2 \\
&= \frac{1}{2^{2n}}\cdot \binom{n}{r}^2
\end{align*}
\begin{align*}
P(X_n = Y_n) &= \sum_{r=0}^n \frac{1}{2^{2n}} \binom{n}{r}^2 \\
&= \frac{1}{2^{2n}}\cdot \binom{2n}{n}
\end{align*}
By symmetry, $P(X_n>Y_n)=P(X_n<Y_n)$
\[
1=P\left(X_n>Y_n\right)+P\left(X_n=Y_n\right)+P\left(X_n<Y_n\right) 
\]
\[
\Rightarrow P\left(X_n>Y_n\right)=\frac{1-P\left(X_n=Y_n\right)}{2}
\]
\[
\Rightarrow P\left(X_n > Y_n\right) = \frac{1}{2^{2n+1}} \cdot \left(2^{2n} - \binom{2n}{n}\right)
\]

So, finally
\begin{align*}
P(X_{n+1}>Y_n)&=\frac{1}{2^{2n+1}} \cdot \left(2^{2n} - \binom{2n}{n}\right)+\frac{1}{2^{2n}}\cdot \binom{2n}{n} \cdot \frac{1}{2}\\
& =\frac{1}{2}
\end{align*}

\newpage

\section{Solution Q5}

Assume that the IDs of the traders are randomly distributed.\\
Let the probability of me winning if I choose my position to be r be F(r).

\[
F(r) = P\left(\text{I win} \mid \text{I am at position } r\right) = P\left(\text{Traders at } 1,\ 2,\ \ldots,\ r-1\ \text{lose}\right) \times P\left(\text{I win at position } r\right)
\]
\[
F(r) =
\begin{cases} 
0 & \text{if } r = 1, \\
\frac{200}{200} \cdot \frac{199}{200} \cdot \ldots \cdot \frac{200-(r-2)}{200} \cdot \frac{r - 1}{200} & \text{if } 2 \leq r \leq 200, \\
0 & \text{if } r > 200.
\end{cases}
\]
Now, suppose my winning probability is maximum at $r=r_0$.\\
Then, $F(r) > F(r-1)$ and $F(r)>F(r+1)$.

\[
\frac{200}{200} \cdots \frac{203-r}{200} \cdot \frac{202-r}{200} \cdot \frac{r - 1}{200} > \frac{200}{200} \cdots \frac{203-r}{200} \cdot \frac{r - 2}{200}
\]
\[
\frac{200}{200} \cdots \frac{202-r}{200} \cdot \frac{r - 1}{200} > \frac{200}{200} \cdots \frac{202-r}{200} \cdot \frac{201-r}{200} \cdot \frac{r}{200}
\]

\vspace{0.5 cm}

\begin{minipage}{0.45\textwidth}
\raggedright
\[
\frac{202-r}{200} \cdot (r-1) > (r-2)
\]
\[
\Rightarrow r^2 - 203r + 202 + 200r - 400 < 0
\]
\[
\Rightarrow r^2 - 3r - 198 < 0
\]
\[
\Rightarrow r \in (-12.65, 15.65)
\]
\end{minipage}
\hfill
\begin{minipage}{0.45\textwidth}
\raggedright
\[
\frac{201-r}{200} \cdot r < r-1
\]
\[
\Rightarrow r^2 - 201r + 200r - 200 > 0
\]
\[
\Rightarrow r^2 - r - 200 > 0
\]
\[
\Rightarrow r \in (-\infty, -13.65) \cup (14.65, \infty)
\]
\end{minipage}

\vspace{0.5 cm}
Hence, by taking the intersection of the above intervals, I conclude that the probability of me winning is maximum at $r=15$.

\end{document}
